\section{Preliminaries: window spaces, Fibonacci evaluation, and Zeckendorf normalization}
\label{sec:preliminaries}

\subsection{Microstate windows and the golden-mean language}
For an integer \(m\ge 1\), let
\[
\Om_m := \set{0,1}^m
\]
be the microstate space of binary words of length \(m\). We also consider the \emph{stable} subset
\[
X_m := \set{x\in \set{0,1}^m : \text{$x$ contains no occurrence of the block \(11\)}},
\]
the length-\(m\) blocks of the golden-mean shift.

\subsection{Fibonacci evaluation}
Let \((\Fib_k)_{k\ge 0}\) be the Fibonacci sequence with \(\Fib_0=0\), \(\Fib_1=1\), and \(\Fib_{k+2}=\Fib_{k+1}+\Fib_k\). For \(\omega=(\omega_1,\dots,\omega_m)\in \Om_m\), define the Fibonacci-weighted evaluation
\begin{equation}\label{eq:Nomega}
N(\omega) := \sum_{k=1}^m \omega_k \,\Fib_{k+1}.
\end{equation}

\begin{remark}[Weight convention]
We adopt the convention that position \(k\) carries weight \(\Fib_{k+1}\). Thus position \(1\) has weight \(\Fib_2=1\) and position \(2\) has weight \(\Fib_3=2\). This is shifted by one from some common numeration conventions; compare \cite{Fraenkel1985,Frougny1992,Frougny1999,Berstel1985,AlloucheShallit2003}.
\end{remark}

\subsection{Zeckendorf normalization and truncation}
Zeckendorf's theorem states that every \(N\in \NN\) admits a unique representation as a sum of non-consecutive Fibonacci numbers \(\Fib_{k+1}\) (\(k\ge 1\)); see \cite{Lekkerkerker1952,Zeckendorf1972,Fraenkel1985}. Equivalently, there exists a unique binary sequence \(z=(z_k)_{k\ge 1}\in \set{0,1}^{\NN}\) such that
\begin{equation}\label{eq:zeck}
z_k z_{k+1}=0\quad\text{for all }k\ge 1,
\qquad
N=\sum_{k\ge 1} z_k \Fib_{k+1}.
\end{equation}
We denote this digit map by \(Z:\NN\to \set{0,1}^{\NN}\), \(Z(N)=z\).

Let \(\proj_m:\set{0,1}^{\NN}\to \set{0,1}^m\) be truncation to the first \(m\) digits,
\(\proj_m(z)=(z_1,\dots,z_m)\).

\subsection{The folding map}
\begin{definition}[Window-level fold]\label{def:fold}
For \(m\ge 1\), define the folding map
\begin{equation}\label{eq:fold}
\Fold_m(\omega) := \proj_m\!\left(Z\!\left(N(\omega)\right)\right)\in X_m,\qquad \omega\in \Om_m.
\end{equation}
\end{definition}

The fold simultaneously performs a \emph{numeric} normalization (passing from a possibly non-Zeckendorf Fibonacci digit string to the unique Zeckendorf digits of the same value) and a \emph{language} normalization (projecting onto the golden-mean constraint). Note that \(X_m\subset \Om_m\) consists precisely of fixed points: if \(x\in X_m\) then \(\Fold_m(x)=x\).

